\chapter{Intrappolate}

«Ahi, mi sono graffiata.»\\
«Di qua, Cate, c’\`e un passaggio.»\\
«Ah, ma perch\'e ho sempre la gonna? Arrivo, Laura.»\\
«Certo che questo giardino non vede un giardiniere da un bel po’ di tempo.»\\
«Giardiniere? Qui ci vorrebbe... non farmelo dire, non sarebbe ecologico. Guarda, l\`i c’\`e un bidone con una rete. Forse ci vive qualcuno?»\\
«Non lo so, questo edificio non \`e neanche sulla mappa. Andiamo a vedere. Rocky dovrebbe essere vicinissimo ormai.»\\
«Il bidone \`e caldo, qualcuno dimora qui. Cosa facciamo, Laura? Rocky potrebbe essere dentro, proviamo a chiamarlo?»\\
«\`E meglio di no. Se c’\`e qualcuno, potremmo mettergli paura, farlo scappare e Rocky con lui. Capiamo se riusciamo a guardare dentro e se c’\`e un punto di accesso sicuro.» Laura si leva lo zaino e prende il drone da rilevamento.
«Cosa fai?»\\
«Cerchiamo un accesso con il drone.» Attiva il drone e fa un volo di prova.
«Ma la porta \`e un po’ aperta, mi sembra.»\\
«Non va bene, ci serve un’entrata che non ci faccia notare. Meglio non rischiare.»

\par\medskip
Caterina si avvicina alla scala antincendio, afferra il corrimano e lo scuote. L'intera struttura vibra e risuona alla sua manovra. Appoggia un piede sul primo gradino. I successivi presentano ferri sporgenti e qualche fessurazione.
«Laura, qui, vieni», le sussurra.

\par\medskip
«Cate, sei gi\`a salita!» Laura spegne il drone e lo riaggancia allo zaino.\\
«Coraggio, seguimi.»

\par\medskip
Caterina avanza per le scale, il tacco ogni tanto si intrappola tra le grate delle scale, ma sembra non curarsene troppo. Laura la segue con passi misurati, senza fare rumore. L'amica si volta verso di lei:
«Hai assolutamente ragione a non entrare nella porta principale. Non stiamo entrando in una proprietà privata» le sussurra con una punta di complicità: «Questo \`e proprio un altro ecosistema, e la prima regola \`e: non perturbarlo» si interrompe: «...ma guarda dove  c’\`e un carrello della spesa!»\\
Con gesto deciso lo spinge via.\\
«Ecco, adesso il passaggio \`e libero» sussurra, divertita. «E senza esserci fatte notare.»

\par\medskip
Laura la raggiunge sul pianerottolo che si piega leggermente: «Non sembra regga bene...»\\
\`E buio ma superano la portafinestra e sono dentro. Laura afferra l'amica per un braccio:
«Fermiamoci un attimo.» Si leva lo zaino, sgancia il drone e dalla tasca laterale sinistra prende il controller.\\

«Cosa vuoi fare?» Caterina le sussurra vicino all'orecchio: «Siamo gi\`a dentro, no?»\\
«Lo attivo comunque. Potrebbe servirci.»\\
\par\medskip
Il LED RX  lampeggia: «Niente connettività, solo registrazione locale. Non basta per la mappa 3D», sussurra.\\
«Qualche problema?»\\
«Non proprio. Avrei voluto costruire una mappa 3D, ma non c'\`e connessione al satellite.»\\


«Wow, Laura! Volevi una mappa di questo posto in 3D? \`E fantastico, ma non sarebbe meglio se cercassimo subito Rocky?»\\
«Hai ragione, forse sono troppo prudente»

\par\medskip
Improvvisamente arriva un latrato poi un urlo, una voce maschile\\
\par\medskip
\begin{center}
  \includegraphics[width=0.7\linewidth]{media/equipaggiamento laura senza sfondo.png}
\end{center}

«Questo \`e Rocky» Attende un istante, poi riprende:  «Sembra venga da sotto. Andiamo a vedere.»\\
«Ma hai sentito quella voce?»\\
«Sì l'ho sentita, adesso sappiamo che non siamo sole di sicuro.» Laura avanza, la stanza \`e buia, lentamente si fanno strada tra una distesa di scranni e  banchi in legno massiccio. La disposizione ricorda quelle di batterie di barricate. Sono  pesanti da spostare, ma non cos\`i fitti da non poter essere aggirati. Raggiungono il corridoio, pochi passi, poi Laura trae tre o quattro respiri veloci per saggiare l'aria: «Odora di limone.» Accende la torcia laser, osserva il soffitto senza illuminare l'ambiente:
«Non ci sono pozzanghere», commenta. «Non \`e acqua infiltrata. Poi il soffitto \`e asciutto, perch\'e il pavimento \`e bagnato?»

\par\medskip
Ancora un latrato. Si dirigono verso le scale. Intorno non vedono nessuno, nessun rumore, tranne quelli provenienti dal piano inferiore.\\
Nessuna luce, tranne i LED RX del ricevitore di Laura.

\par\medskip
«Scendiamo, Cate. Attenta...» le dice, le scale non sono tanto sicure. L'ultima rampa le conduce in una grande hall. Da una stanza proviene una luce tremolante, che le permette di distinguere alcuni cartoni a terra e delle strutture in legno lungo il perimetro principale. Laura cataloga tutto nel cervello, ma la costruzione della mappa mentale \`e interrotta da un altro latrato. Il suono non \`e lontano, ma arriva sordo, rimbombante. Accanto alla tromba delle scale ci sono due porte, una per lato. Quella a destra: da lì proviene il suono.\\
Entrano: sono i bagni del convitto. Un lavello comune e tre ritirate, due hanno le porte ridotte come le altre a nudi passaggi la terza \`e ancora integra e chiusa. Il suono proviene dal basso, forse da una delle turche. Caterina apre la porta ed entra in uno dei gabinetti, Laura ne ispeziona un altro, ma l\`i non percepisce alcun rumore.
«Vieni a sentire» la chiama Caterina con un filo di voce: «da qui, arriva un respiro.»

\par\medskip
«Fammi sentire» si abbassa ai piedi dell'amica ma in un attimo… \emph{sbam!}
La porta si chiude e scatta una serratura.

\par\medskip
Caterina grida. Qualcuno le ha chiuse dentro.
